\documentclass{jsarticle}
\usepackage[dvipdfmx,final]{graphicx}
\usepackage{amsmath}
\usepackage{amssymb}
\begin{document}
\title{}
\author{}
\date{}
\maketitle
\thispagestyle{empty}
\includegraphics[width=8cm,clip]{M1100010}

ユウの燃え上がる目

\includegraphics[width=8cm,clip]{M1100032}

私に対しての反抗期


昨年は、大変お世話になりました。新年あけましておめでとうございます。
今年もよろしくお願い申し上げます。
２０２１年は、ユウくんの写真の誤解を招くことをしてしまい、
申し訳ございませんでした。２０２２年にユウくんが
あっかんべーを２０２１年のはがきを出したあとに、またやってくれて、
２０２０年と２０１９年ずっと、あっかんべーしてくれていて、
恥ずかしくて、ビデオに撮っていたのを消しまくり、未練がましく、２０２１年に
似たような写真を載っけて、フォトリーディングとSRSのしているユウの写真が
邪魔だったかと、この子瞑想法もするので、これももったいないけど、
消しまくり、この出したあっかんべーも今では、消しています。
２０２２年の写真、大切にしてね。あっかんべーでない写真を誤解されたのかと
気づいたときに、私、殺菌されていました。家族ですね。
あと、医王寺のお経の本と四国巡礼のお経の本の、無意識という単語は、
凝視は、書き位置が問題だったので、気づいてくれるかなとおもっていたが、
一向に気づいてくれずに、医王寺は、私が購入した存在と無のサルトルの本の
アドバイスとおもうし、今は、この本の書き位置は大丈夫になっている。
ご先祖様は、私の作曲できる、音を一切使わずに作る曲の能力は、
姉が、HANONで音を使わないように弾きなさいと、姉の視力で、
本や楽譜を風景で速読をモネの睡蓮の絵画と同じ見方の、音を使わずに、
その上に、弾く曲がどの音階を使っているかで、楽譜を見ながら、
鍵盤を見ずに弾く能力にべたぼれして、２００７年にご先祖様が、
速読いいと、私らがいるから、速読やりなさいと言っているしか、
捉えようない出来事と思い、ちゃんと父方の先祖供養に来るべきともおもいます。


\end{document}
